\ficha{Ferguson e Schularick (2012), The Thin Film of Gold}{fergusonschularick2012}

\begin{identificacao}
FERGUSON, Niall; SCHULARICK, Moritz. The ``thin film of gold'': monetary rules
and policy credibility. \textit{European Review of Economic History}, v. 16,
n. 4, p. 384-407, 2012. DOI: 10.1093/ereh/hes006.

Artigo de história econômica quantitativa, em inglês. Lido integralmente na
versão de working paper, NBER Working Paper 13918, abril de 2008, cujo título
traz \textit{in Developing Countries}, 38 páginas. As páginas citadas abaixo
remetem à paginação do working paper, não à do artigo publicado.
\end{identificacao}

\bloco{Problema e tese}
O texto pergunta se um país em desenvolvimento ganha credibilidade ao submeter a
política monetária a uma regra estrita, e trata o padrão-ouro de 1880-1914 como
experimento natural. A resposta é negativa para a periferia pobre. O prêmio de
risco cai após a adesão ao ouro nos países já desenvolvidos, mas não nos pobres,
onde o investidor precificava vulnerabilidade a choques econômicos e políticos, e
não inconsistência temporal. O ouro não era selo, era película fina, e o mercado
olhava atrás dela.

\bloco{Argumento}
\begin{enumerate}[leftmargin=*,nosep]
\item A hipótese sob teste é a de Bordo e Kydland e Bordo e Rockoff: o padrão-ouro
operou como regra com cláusulas contingentes de escape, e o mercado cobrava menos
de quem a adotava, 40 pontos base a menos num título em ouro (p. 6). Obstfeld e
Taylor estimam até 30 pontos base e sustentam que o selo prevalecia sobre a
dívida pública e os termos de troca (p. 6).
\item A divergência da literatura é atribuída a três causas: a codificação do que
é estar no ouro; a especificação, em que a variável do ouro vira proxy de
omitidas, como no Japão de 1897, ano de adesão mas também de vitória militar e de
conversão da dívida (p. 7-8); e, sobretudo, a amostragem, pois Bordo e Rockoff
usam nove países, Obstfeld e Taylor 17 e Flandreau e Zumer 21, quando mais de
sessenta tomadores tinham títulos em moeda forte listados em Londres (p. 8).
\item Restringindo a base nova às amostras antigas, os resultados antigos se
reproduzem, com queda de 30 a 40 pontos base (p. 15-16). A discórdia com Bordo e
Rockoff não é de técnica, é de população de países.
\item Na amostra cheia, de 58 tomadores, o coeficiente do ouro cai a cerca de oito
pontos base e não chega perto da significância (p. 18). Nada disso muda com a
codificação de facto, com a matriz de Meissner, com período de prova de cinco
anos, nem restringindo a amostra a quem veio do papel para o ouro (p. 20-21).
\item A quebra por subamostra localiza o efeito. Nas 23 colônias britânicas o
regime monetário é irrelevante (p. 22). Entre os independentes desenvolvidos o
ouro vale até 50 pontos base, significativo a 1\%, como em Bordo e Rockoff. Entre
os 22 independentes pobres o coeficiente tem sinal errado e é insignificante em
toda especificação (p. 25). A explicação é dupla: economias pobres eram mais
expostas a choques agrícolas, e nelas o interesse agrário pesava mais do que
banqueiros e rentistas, o que fragilizava a promessa de autolimitação (p. 28-29);
e, onde a regra política é reescrita com frequência, o regime monetário vira
preocupação de segunda ordem (p. 30).
\end{enumerate}

\chave{The market, we infer, did not confer a `good housekeeping seal of approval'
on poor peripheral countries merely because they adopted the gold standard.}{p. 25}

\bloco{Conceitos e definições operacionais}
Regra monetária é pré-compromisso que amarra as mãos de quem governa contra o uso
da política monetária para ganho de curto prazo; o padrão-ouro é o caso histórico
e o currency board o equivalente contemporâneo (p. 2). Credibilidade é
operacionalizada como prêmio de risco-país, o diferencial entre o rendimento do
título soberano em ouro ou em libras e o do consol britânico (p. 14). Estar no
ouro é codificado de jure, com conversibilidade legislada e mantida na prática, e
testado contra a codificação de facto e a matriz de Meissner (p. 12). País em
desenvolvimento é o de renda per capita abaixo de um terço da britânica, cerca de
1.500 dólares de 1990 em 1900 (p. 24, nota 21). Período de prova é a variável que
só liga após cinco anos de cumprimento da regra (p. 13).

\bloco{Evidência e método}
Painel anual de 1880 a 1913, 34 países independentes e 23 colônias britânicas,
1.449 observações de prêmio, cerca de três vezes a maior base anterior.
Rendimentos coletados à mão em \textit{The Investor's Monthly Manual} e
\textit{The London Stock Exchange Weekly Intelligence}; controles no
\textit{Statesman's Yearbook}, no \textit{Fenn's Compendium} e nos relatórios da
Corporation of Foreign Bondholders (p. 9). Duas baterias de controles, uma
histórica e outra moderna, dão resultados próximos (p. 17). Estimação por
Prais-Winsten com efeitos fixos e erros padrão corrigidos para painel, betas por
país sobre um spread mundial, e robustez em GMM sistêmico de Arellano e Bover.
Testes de Chow justificam separar colônias, ricos e pobres. A tabela VI mostra a
periferia com metade da renda per capita e da abertura comercial da amostra
avançada, e 3,88 anos de conflito contra 0,61 (p. 30).

\bloco{Agentes e vertentes}
Pela regra: Bordo, Kydland, Rockoff, Obstfeld, Taylor e Meissner. Contra:
Flandreau e Zumer, que põem serviço da dívida, crescimento e inflação no lugar do
regime, além de Eichengreen e Hausmann, Mitchener e Weidenmier e os próprios
autores em 2006, com o efeito império. Na teoria, Kydland e Prescott e Barro e
Gordon para inconsistência temporal, e Drazen e Masson para a distinção entre
credibilidade das políticas e credibilidade de quem as formula. Os agentes
históricos aqui são grupos, não pessoas: o interesse agrário contra banqueiros e
rentistas (p. 29), e a imprensa financeira da City, que escrevia sobre a política
desses países e quase nunca sobre conversibilidade (p. 31).

\bloco{Críticas e limites}
Os autores admitem a endogeneidade da adesão, de modo que o efeito estimado é
média parcialmente incondicional dos países em condição de adotar a regra (p. 8,
nota 8), e admitem que o coeficiente achado nos ricos pode ser proxy de melhoras
não captadas ou efeito de custo de transação (p. 31). O objeto medido é estreito:
prêmio soberano em Londres, não volume de crédito, prazo, juro interno nem
câmbio. Duas inconsistências internas do working paper pedem cuidado antes de
citar número: a tabela II rotula a amostra de Bordo e Rockoff como ``BR (1998)''
e traz sete países, enquanto o texto fala em nove (p. 8, p. 16); e o texto fala em
16 desenvolvidos independentes, enquanto a tabela IV traz 13 (p. 23-24).

O Brasil não é objeto do texto. Aparece uma vez, como país com estoque de
investimento estrangeiro em torno de 100\% do PIB em 1913 (p. 5). Não há menção à
Caixa de Conversão, ao Convênio de Taubaté nem à adesão de 1906, e o painel
termina em 1913, de sorte que o Brasil no ouro entra tarde e diluído numa média.

\chave{For a poor country seeking to borrow in London at sustainable rates, we are
tempted to suggest, it made more sense to become a British colony than to join the
gold standard.}{p. 28}

\begin{dialogo}
\item Procurar nos editoriais favoráveis à Caixa a promessa de que a estabilização
baratearia o crédito brasileiro em Londres. Vocabulário: ``credito do paiz'',
``confiança dos capitaes estrangeiros'', ``praça de Londres''. Se aparecer forte e
sem qualificação, os defensores operavam com a hipótese de Bordo e Rockoff, que
este texto nega para país pobre, o que reclassifica o argumento de 1906 como
crença dos agentes, não como mecanismo comprovado.
\item Procurar a contraparte cética, a alegação de que o investidor olha o
orçamento e a ordem política antes do regime monetário. Vocabulário:
``equilibrio orçamentario'', ``serviço da divida'', ``renda das alfandegas''.
Testa se a leitura de Flandreau e Zumer, aqui endossada, tinha equivalente
doméstico, e de que campo ela vinha.
\item Verificar o que os diários transcreviam da imprensa financeira britânica
sobre o Brasil. Os autores dizem que ela tratava da periferia pela política e
quase nunca pela conversibilidade (p. 31). O que se reproduzia mede a distância
entre a repercussão suposta em Londres e a real.
\item Procurar os grupos do modelo nas colunas brasileiras: lavoura e exportadores
contra praça bancária e portadores de apólices. Testa a proposição de que na
periferia o interesse agrário superava o dos credores, hipótese que ajuda a
explicar por que o eixo do debate de 1906 foi o nível da taxa, e não a existência
da regra.
\item A suspensão de 1914 testa a própria tese. Procurar, antes de 1914,
articulista que antecipe que a Caixa cederá ao primeiro choque externo. Se
houver, a expectativa de que a regra não se sustentaria também era doméstica.
\end{dialogo}

\usonocapitulo{Parte I, como contraponto empírico direto a Bordo e Rockoff (1996).
Sustenta duas afirmações. Primeira: o selo de credibilidade, medido pelo prêmio de
risco soberano, não se verifica para a periferia pobre, o que retira da adesão
brasileira ao ouro a justificativa mais forte que a literatura lhe deu. Segunda: a
assimetria centro-periferia aqui é resultado medido, com corte de renda e teste de
Chow, e não postulado teórico, o que a torna utilizável ao lado de Franco (1988) e
de Marcondes. O texto não trata do Brasil nem da Caixa de Conversão, e toda ponte
com as fontes primárias é hipótese a testar, não referência a citar.

\chave{it may make more sense to think of the gold standard less as a `seal of
approval' and more as a kind of `thin film', behind which investors were wise to
look.}{p. 32}}
